% RESEARCH_DRAFT_SHA256: 1d1891d37d6543968d061333525c2de76314f8846370d497dfa5ba10d56c064a
% Layout and source-reference wrappers expanded for the web renderer.
\documentclass[11pt]{article}
\usepackage[margin=25mm]{geometry}
\usepackage{fontspec,unicode-math}
\setmainfont{Latin Modern Roman}
\setmathfont{Latin Modern Math}
\usepackage{amsmath,amssymb,mathtools,xurl,hyperref}
\hypersetup{colorlinks=true,urlcolor=blue,linkcolor=blue}
\setcounter{secnumdepth}{0}
\setlength{\parindent}{0pt}
\setlength{\parskip}{5pt}
\setlength{\emergencystretch}{3em}
\newcommand{\sref}[2]{\hyperlink{source:#1:#2}{[S#2]}}
\newcommand{\cataloguescope}[1]{\paragraph{Source record.}#1}
\newcommand{\floor}[1]{\left\lfloor#1\right\rfloor}
\begin{document}
\section*{Research-attempt provenance}
This standalone research attempt retains the complete original definition body below,
including its historical source and preparation notes. Those notes describe the
earlier catalogue preparation; the new attempt and its verification are identified
separately after that body. No original ID, catalogue status, or source statement is changed.

\section{5.8 (Schleimer) — Why SnapPy works in practice}\label{statement}

\subsection{Definitions and mathematical statement}

\paragraph{Shared notation.}Standard number systems are $\mathbb N,\mathbb Z,\mathbb Q,\mathbb R,\mathbb C$. Specialized notation and index ranges follow the cited source and the definitions below.


A triangulation describes a space by simplices and face identifications; an ideal triangulation places designated vertices at ends or cusps. In hyperbolic $3$-manifold computations, a tetrahedron can be described by a complex shape parameter, and gluing equations impose compatibility around edges and peripheral curves. Numerical procedures approximate solutions to these equations; a certificate requires a specified exact or interval-based verification condition. SnapPy is software working with such triangulated $3$-manifolds and related topological/geometric data.
\paragraph{Record scope and qualifications.}
The catalogue title is retained as metadata. Named certificates, contracts, bounds or frameworks in that title are not asserted as results of the original source extract.
The 2015 list attributes this agenda to Schleimer. It does not specify a SnapPy algorithm, input distribution, success metric or theorem. The context definitions describe terminology from the official documentation; no guarantee for all SnapPy outputs or explanation of its observed performance is claimed.
\paragraph{Statement recorded in the supplied source.}
Give a rigorous explanation for why SnapPy works so well in practice. \href{https://arxiv.org/pdf/1512.04620#page=4}{[S2]}
\paragraph{Source status.}
The supplied catalogue labels this source record open; that metadata was checked against the supplied snapshot on 2026-10-11. The label is preserved for provenance and is not an independent verification of present open status or mathematical correctness.
\subsection{Short English statement}
Explain, in a rigorous and explicitly chosen computational model, the practical effectiveness of SnapPy.
\subsection{Sources}
\begin{itemize}
\item[S1] Ulam AI and the UnsolvedMath Contributors, UnsolvedMath: A Curated Collection of Open Mathematics Problems, catalogue entry. Catalogue curation and metadata: CC BY 4.0. The source extract is retained; definitions, English summaries and qualifications are editorial additions.
\url{https://huggingface.co/datasets/ulamai/UnsolvedMath}.
\item[S2] AMR Open Problem Lists, source linked by the catalogue.
\url{https://arxiv.org/pdf/1512.04620#page=4}.
\item[S3] SnapPy official documentation, Triangulation: triangulations and gluing equations. Consulted for terminology, not a universal performance theorem.
\url{https://snappy.computop.org/triangulation.html}.
\end{itemize}
\paragraph{Preparation and review.}
The original source and UnsolvedMath contributors retain attribution for the source extract. Codex (GPT-6; exact serving revision not exposed) assisted with definitions, English summaries, scope notes and formatting. No new human mathematical authorship or independently verified solution is asserted. The supplied catalogue credits: Schleimer. Catalogue identity checked against the complete supplied registry and current repository. Distinct upstream records, including repeated questions and broad or damaged extracts, are retained for complete coverage of the source CS category. A repeated record is not a new independent conjecture. Complete mathematical-equivalence and current-literature review has not been performed. Preparation and final review follow contribution prompts C and D, with source gaps disclosed. Structural and rendering checks do not establish a new theorem or open status.
\label{end:statement}

\section{Literature and context check}
The original list's Problem 5.8 and SnapPy's official verified-computation
documentation were consulted. The documentation describes interval certificates
and explicit failure behavior. These concern correctness when a check succeeds;
they do not by themselves explain a broad practical success rate. The following
standard contraction argument is reproved rather than attributed as new theory.
\section{Research attempt: A quantitative local certificate and the missing global explanation}
Let $F$ be continuously differentiable on the closed ball $B(z_0,r)$ in a
finite-dimensional real or complex vector space, and let $A$ be an invertible
linear map. Suppose
\[
 \beta=\|AF(z_0)\|,\quad
 \sup_{z\in B(z_0,r)}\|I-ADF(z)\|\le\kappa<1,
 \quad \beta\le(1-\kappa)r.
\]
For $T(z)=z-AF(z)$, integrating its derivative along a straight segment gives
$\|T(z)-T(w)\|\le\kappa\|z-w\|$. Also
$\|T(z)-z_0\|\le\kappa r+\beta\le r$. Starting from $z_0$, the iterates stay
in the ball; consecutive differences shrink geometrically, so the sequence
is Cauchy. Its limit $z_*$ belongs to the closed ball and satisfies
$T(z_*)=z_*$ by continuity, hence $F(z_*)=0$ because $A$ is invertible.
Two fixed points would have distance at most $\kappa$ times their distance,
so the zero in the ball is unique. This proves the certificate in full.

As a concrete complex calculation take $F(z)=z^2-z+1$,
$z_0=1/2+7i/8$, $r=1/64$, and $A=(2z_0-1)^{-1}=-4i/7$.
Then $F(z_0)=-1/64$, $\beta=1/112$, and
\[
 |1-AF'(z)|=2|A|\,|z-z_0|\le1/56=\kappa.
\]
The self-map inequality is $1/112\le55/3584$, which is true. The disk lies
strictly in the upper half-plane since $7/8-1/64>0$; it therefore certifies a
positive-imaginary root. This is an exact rational certificate for a shape-type
polynomial, not a check of all equations of any claimed manifold.

The failed shortcut would infer certification from a small numerical residual
alone. For $F_\varepsilon(z)=\varepsilon(z-1)$ and $z_0=0$, the residual tends
to zero with $\varepsilon$, while the zero remains at distance one. The inverse
derivative has norm $1/\varepsilon$, and $\beta=1$ correctly exposes the
problem. A residual tolerance that ignores inverse-Jacobian scale is therefore
not a proof.

Applied to a selected square gluing system, the certificate still requires
proof that the selected equations imply the remaining edge and peripheral
conditions, that the ball stays in the appropriate shape domain, and that the
geometric interpretation follows. Explaining practice additionally requires a
reason why the numerical search finds such well-conditioned balls on a specified
input class. None of these global or statistical facts follows from the local
contraction theorem.
\section{Verification and adversarial audit}
Every constant in the scalar example is exact, and the upper-half-plane
condition was checked on the whole disk. The existence proof uses convergence
of the iterates rather than treating a numerical iterate as an exact root.
SnapPy itself was not run; no census-wide or manifold-specific certificate is
claimed.
\paragraph{Audit corrections.} A local polynomial certificate was kept distinct from a manifold hyperbolicity proof and from an explanation of practical performance; residual size alone was falsified.
\paragraph{External computation.} None; symbolic arguments and boundary checks only.
\section{Scope and remaining gaps}
The original question is an agenda without one algorithm, distribution or
success criterion. The first analytical gap for the explored route is a useful
bound on finding a certificate, including condition numbers and appropriate
initialization, for an explicitly stated family of actual manifold inputs.
\paragraph{Integration review.} Codex (GPT-6; exact serving revision not exposed) reviewed the written arguments and expanded layout and source-reference wrappers. Repository metadata and the production web renderer were checked using Python 3.13 and JavaScript with MathJax 3.2.2. This model-assisted review is not independent human proof verification; no PDF compilation or formal proof checking is asserted.

\section{Completion Estimate}
\noindent\textbf{COMPLETION: 0\%}
The calculations below the preserved source statement investigate a concrete route,
but establish no substantiated new progress toward the entire original target.
Known special cases and failed sufficient conditions are not credited as a solution.
\section{References}
The original source references and catalogue attribution are preserved in the
Sources section above. The following URLs were actually consulted for this attempt;
retention of an additional original reference does not claim fresh consultation of it.
\begin{itemize}
\item \url{https://arxiv.org/pdf/1512.04620} (consulted 2026-10-11). 
\item \url{https://snappy.computop.org/verify.html} (consulted 2026-10-11). 
\end{itemize}
\section{Model and contributor disclosure}
OpenAI gpt-6-astra, requested reasoning configuration ultra, generated this research
attempt on 2026-10-11. The exact serving revision is not exposed. Codex GPT-6
performed orchestration and repository-format review; its exact serving revision
is also not exposed. The mathematical route, boundary checks and final adversarial
self-audit followed contribution prompts A then B in
\texttt{docs/contribution\_prompts.md}, after reading \texttt{CONTRIBUTING.md}.
Python 3.14.2 assembled the standalone documents; any mathematical computation is
specified in Verification. No theorem prover was used. No human mathematical
authorship, substantial human proof correction, or independent human proof checking
is asserted. Model-assisted review, source retrieval, and structural checks do not
establish correctness or novelty. Full private reasoning transcripts are unavailable
as reproducibility artifacts; the complete arguments, inputs and outcomes actually
used are stated here or in the companion batch audit.
\section{Final status}
Unresolved. This is a research attempt with an explicit failed route or restricted
calculation, not a full solution, not a current-status certification, and not a claim
that the elementary material is new. The catalogue's original source-status label
is preserved independently of this attempt status.
\end{document}
